%! Author = lili
%! Date = 6/9/26

% \chapter{BW4 - Chromosomentheorie, Genkopplung, genet. Rekmb}

% bis F9 ist Wdh.

%\section{Chromosomentheorie}
\defin{zweichrchr}
\begin{figure}[H]
    % TODO
    \caption{\gls{zweichrchr}}
\end{figure}

pro Chromatid [bzw. Chromosom] gibt es einen durchgehenden, linearen DNA-Doppelstrang
der zwei ``offene'' Enden hat; diese Enden sind durch Telomere speziell geschützt

% todo ref def:semiRep

% TODO Das Problem mit dem Ende ... F 12 

\section{Thomas Hunt Morgan und die Chromosomentheorie der Vererbung}
Prägte die Begriffe \gls{Mutante} vs. \gls{Wildtyp}  und\gls{gen} vs. \gls{Allel}.
Vater des \gls{cm}. % TODO cm F 41
% TODO F 15-16  bsp warum guter Modelorganismus

\subsection{Krezungsversuche mit Drosophila melanogaster}
% TODO F 18-23 Experiment

\defin{autosom}
\defin{heterosom}
% TODO unterschied zu chromosom? F 23

% TODO F 25-31 Experiment

\subsection{Rekombinationsfrequenz}
\defin{rekombinationsfrequenz}
% TODO F 32 - 33

\subsection{Kopplungskarten}
% TODO F 34 -35

\subsection{Schlußfolgerungen}
% TODO als Karteikarte
\begin{itemize}
    \item Gene sind linear auf den Chromosomen angeordnet
    \item Je weiter Gene auseinander liegen, desto
    wahrscheinlicher ist ein Austausch (Crossover)...
    \item ... je näher sie zusammen liegen, desto
    unwahrscheinlicher
    \item Austauschhäufigkeiten stellen ein Maß für die relative
    Lage der Gene zueinander dar
    \item Kopplungsgruppen
    \item Genetische Karten / Kartierung von Chromosomen
\end{itemize}

\subsection{Zusammenfassung: Unterschiede zu Mendel}
% TODO F ?? - 36

\section{Genetische Kartierung}
% TODO F 39 -40, 43 - 45

% TODO frage Unterschied F 42

\section{Erbkrankheiten bei H. sapiens}
\begin{bsp}
    \textbf{Hämophilie}
    % TODO F 46 - 48
\end{bsp}
\subsection{multifaktorielle Erkrankungen}
% TODO defi usw F 49
\subsection{Diagnose}
% TODO F 50

